\documentclass[authoryear,preprint,review,12pt]{elsarticle}


\textheight24cm
\textwidth15.5cm              
\oddsidemargin0.25cm
\topmargin-1.8cm
\usepackage[utf8]{inputenc}
\usepackage{booktabs} % For professional looking tables
\usepackage{tabularx} % For tables with adjustable-width columns
\usepackage{ltablex} % Keeps tabularx environment with added features
\usepackage{geometry} % For setting page margins
\usepackage{caption} % For captions
\usepackage{amsmath}
\usepackage{tcolorbox}  % 高亮和文本框支持
\usepackage{soul}


\usepackage{caption}
% Adjust the page margins to fit the table size
\geometry{
  left=0.75in,
  right=0.75in,
  top=1in,
  bottom=1in
}


\usepackage[T1]{fontenc}
\usepackage{lmodern}

\usepackage{xeCJK}


%Packages
\usepackage{pifont}
\usepackage{placeins}
\usepackage{color}
\newcommand{\red}[1]{\textcolor{red}{#1}}
\usepackage{graphicx}
\usepackage{pdfpages}
\usepackage{lineno}  
\usepackage{subcaption} %
\usepackage{amsmath} % 引入必要的包
\usepackage{amssymb}
\usepackage[utf8]{inputenc}
\usepackage{float}
\usepackage{arydshln}
\usepackage{url}
\def\UrlBreaks{\do\/\do-}
\usepackage{multirow}
\usepackage{comment}
\usepackage{threeparttable}
\usepackage{bm}
\usepackage{amsmath}
\usepackage{booktabs} % For professional looking tables
\usepackage{tabularx} % For tables with adjustable-width columns
\usepackage{geometry} % For setting page margins

%\usepackage[x11names]{xcolor}
\usepackage[
  colorlinks,
]{hyperref}
\definecolor{DodgerBlue4}{rgb}{0.0627, 0.3059, 0.5451}
\AtBeginDocument{\hypersetup{citecolor=DodgerBlue4}}

%%%%%%%%%%%%%%%%%%%%%%%%%%%%%%%%%%%%%%%%%%%%%%%%%%%
\journal{International Journal of Geographical Information Science}
\begin{document}
\begin{frontmatter}
\title{Measuring causal associations by geographical cross mapping cardinality}

\author[a,b,c]{Wenbo Lv}
\ead{wenbo.lv@snnu.edu.cn}

\author[d]{Yongze Song}
\ead{yongze.song@curtin.edu.au}

\author[c,e]{Wen Yi}
\ead{wen.yi@polyu.edu.hk}

\author[f]{Shaoqing Dai}
\ead{shaoqing.dai@whu.edu.cn}

\author[b]{Wufan Zhao\corref{cor1}}
\cortext[cor1]{Corresponding author: Wufan Zhao, wufanzhao@hkust-gz.edu.cn}
\ead{wufanzhao@hkust-gz.edu.cn}


\address[a]{School of Geography and Tourism, Shaanxi Normal University, Xi'an, China}
\address[b]{Urban Governance and Design Thrust, Society Hub, The Hong Kong University of Science and Technology(Guangzhou), Guangzhou, China}
\address[c]{Department of Building and Real Estate, The Hong Kong Polytechnic University, Hong Kong, China}
\address[d]{School of Design and the Built Environment, Curtin University, Perth, Australia}
\address[e]{The Hong Kong Polytechnic University, Shenzhen Research Institute, Shenzhen, China}
\address[f]{School of Resource and Environmental Sciences, Wuhan University, Wuhan, China}

%%%%%%%%%%%%%%%%%%%%%%%%%%%%%%%%%%%%%%%%%%%%%%%%%%%
\begin{abstract}


Contents for abstract


\end{abstract}

\begin{keyword}
Causal associations \sep spatial empirical dynamic modeling \sep spatial cross sectional data \sep intersection cardinality \sep geographical cross mapping cardinality
\end{keyword}
\end{frontmatter}

%%%%%%%%%%%%%%%%%%%%%%%%%%%%%%%%%%%%%%%%%%%%%%%%%%%
\linenumbers


\section{Introduction}

Contents for Introduction

\section{Methodology} % or Material and methods

Contents for Methodology

\subsection{Phase space reconstruction}

\subsection{Intersection cardinality}

\subsection{Causal scores}

\section{Results}

\subsection{Case1}

\subsection{Case2}

\subsection{Case3}


\section{Discussion}

Contents for Discussion

\section{Conclusion}

Contents for conclusion,likes:
This study offers a new method for understanding traffic accident patterns in Germany by utilizing continuous, detailed data and advanced geospatial modeling. We developed the RGOZH model by introducing numerical optimization into GOZH to make the analysis more reliable. The results of this study provide solid quantitative evidence for policies and behaviors related to regional development inequalities in Germany. On the one hand, the differences in the German federal government's various regional subsidy policies largely align with the current pattern of traffic accidents in those regions. On the other hand, based on the results of this study, a more rational traffic accident regulation program can be developed in the future. This includes designing different regulatory approaches, policies, actions, and budget allocations for different regions and demographic groups, such as the East-West Germany boundary, as identified in this study. The main motivation of this paper is to propose a new GOZH-based approach to study quantifiable outcomes of geographic heterogeneity, and we have chosen the more representative German traffic accident study as a case study.

\section*{Disclosure statement}

No conflict of interest exists in this manuscript, and the manuscript was approved by all authors for publication.

\section*{Funding}

\section*{Data availability statement}

All the data used in this study are available from the website (https://regionalatlas.statistikportal.de/) of Federal Statistical Office of Germany.


%\section*{References}
\baselineskip12pt
\bibliographystyle{elsarticle-harv} 
\bibliography{references.bib} 


\end{document}
